\section{Discussion}\label{sec:discussion}

\paragraph{What the framework makes possible}
The usual argument about antisemitic incidents asks whose count is right.
The framework replaces that question with a different one: what is the
population that every list samples, and how does each list sample it? Once
the target is defined (\cref{sec:framework}), the police count, an NGO audit
and a community-security log become \emph{data about} the same unknown
total, not competitors for it. Their overlaps are informative, their
dependence is modelled rather than assumed away (\cref{prop:referral}), and
the part of the answer that the data cannot deliver is isolated in a single
interpretable parameter $\omega$ (\cref{thm:closed,cor:bounds}). The
sensitivity frontier tells a reader how much unmeasured dependence a stated
conclusion can withstand (\cref{thm:id-frontier}). Linkage
error, which in practice is often handled silently, becomes an explicit
operator whose effect on the estimate can be quantified and corrected
(\cref{thm:linkbias,cor:linkcorrect}). Its residual uncertainty is carried
into the final interval (\cref{alg:pipeline}).

The funnel decomposition (\cref{prop:funnel,thm:shapley}) turns an estimate
of the dark figure into something a policymaker can act on. A gap between
two jurisdictions' recorded rates can arise because victims there are less
willing to come forward, because the police are less likely to recognise a
bias motivation, or because the agency does not submit its data. Each of
these calls for a different response: community outreach, officer training
or reporting infrastructure. Raw counts cannot tell these explanations apart.
The decomposition can, to the extent that the stage-specific rates are
identified by the auxiliary data of \cref{tab:funnel-sources}. Incidents that
were never disclosed to anyone remain beyond the reach of every list, so
that stage rests on victim surveys.

\paragraph{Limitations}
We stress five.

\emph{Untestable assumptions.} No method escapes the basic
non-identifiability of the unobserved cell (\cref{thm:nonid}). Our approach
makes the identifying assumption explicit and adjustable, but it does not
remove it. A sensitivity interval for $\omega$ is only as credible as the
range chosen for it, and that range must be argued from knowledge of how the
lists are generated, not from the data. It must in general be two-sided:
the model without a $K$-way interaction can overstate the dark figure as well
as understate it (\cref{prop:direction}, \cref{ex:id-hetero}). A one-sided
range is justified only by a mechanism of known sign, such as referral
between two lists (\cref{rem:id-calibrate}); with three or more lists,
referral itself can push the bias either way. Latent-class models replace the
assumption of no highest-order interaction with conditional independence
within classes. That is weaker in some directions and not in others. The
Chao bound is a lower bound only when heterogeneity runs across incidents.
Heterogeneity across lists can turn it into an upper bound
(\cref{prop:id-mt}), so it is a diagnostic, not a floor.

\emph{Scope harmonisation.} MSE requires every list to be sampling the same
population. The FBI records criminal offences with a recorded bias
motivation, whereas NGO audits also record non-criminal incidents under
their own and evolving criteria \citep{adl2025audit}. Harmonisation rules
(\cref{sec:src-harmonise,sec:plan}) can align the lists only as far as
their records contain the information needed to apply a common
definition. Residual mismatch
behaves like coverage error and cannot be fully separated from it.

\emph{Data access.} The methods assume record-level access to several
lists. Police records, NGO case files and community reports contain
sensitive information about victims who did not consent to being linked.
Obtaining such data will require data-sharing agreements,
privacy-preserving linkage \citep[e.g.,][]{schnell2009} and disclosure
control for small cells \citep{hundepool2012}, possibly with formal
privacy guarantees \citep{dwork2014}. Each of these degrades linkage
quality or precision to some degree. Encodings for privacy-preserving
linkage can themselves be attacked (\cref{sec:ethics-pprl}), and the ethics
protocol of \cref{sec:ethics} may rightly rule out some linkages
altogether.

\emph{The closed-population idealisation.} Incidents occur in continuous
time, lists have different reporting windows and lags, and records are
added or revised after the reference period ends. Treating each state--year
domain as a closed population is a convenient approximation. Its error is
concentrated at period boundaries and in lists with long reporting delays.

\emph{Approximate linkage-error models.} The operator
$\Lop(\delta,\varepsilon)$ assumes that detachments and false merges occur
independently across records, at rates that are at most list-specific
(\cref{rem:link-listspecific}). Real linkage errors
cluster: they are more frequent for incidents in dense places, for generic
incident descriptions and for records with missing fields, and they may be
correlated with capture history itself. The correction of
\cref{cor:linkcorrect} should therefore be read as a first-order adjustment.
Propagating linkage uncertainty by mixing over linkage draws
(\cref{alg:pipeline}) is preferable whenever record-level comparison data are
available. The modular posterior used there (\cref{eq:models-cut})
deliberately ignores the information that the capture model carries about
which records match, trading some efficiency for robustness.

\paragraph{Extensions}
Three extensions seem most valuable. First, \emph{open populations and
continuous time}: modelling incident occurrence as a point process and each
list as a time-varying, delayed thinning of it would remove the
closed-population idealisation and connect dark-figure estimation to
nowcasting of reporting delays. Second, \emph{online incidents}: here the
unit of analysis is itself unsettled, and platform moderation and automated
detection act as lists with their own selection and classification error, a
direct analogue of the funnel's classification stage. Third,
\emph{international comparison}. In the United Kingdom the Community Security Trust maintains a
long-running list of antisemitic incidents and documents the share it receives
from the police \citep{cst2026incidents}. The European Union Agency for
Fundamental Rights has surveyed Jewish people in several member states about
their experiences of antisemitism and about non-reporting
\citep{fra2018antisemitism,fra2024antisemitism}; see
\cref{sec:src-community,sec:src-surveys}. Combining such sources country by
country, with the funnel decomposition as a common accounting framework, would allow cross-national
comparisons of reporting, not only of recorded counts.

\paragraph{Why measure what is missing}
A count that omits its dark figure can mislead in both directions. Rising
counts may reflect greater willingness to report, better classification or
wider agency participation as much as rising incidence. Falling counts may
reflect lost trust or lapsed reporting. Resources allocated in proportion to
recorded incidents reward the communities and agencies that report, which
are not necessarily those that are most affected. Nor is there reason to
expect the dark figure to be random noise. If it is concentrated where victims
expect least from reporting, its size is itself a measure of the distance
between a community and the institutions meant to protect it. Estimating
it honestly, with explicit assumptions, propagated uncertainty and a plan fixed before the data are
seen, is how a measurement programme takes seriously the incidents that
nobody wrote down.
